\section{Curvature, torsion, and Einstein--Cartan--Holst action}
\label{chap:gravedad-cartan-holst}

The HMT transport structure distinguishes the memory variable,
holonomy, and closure defect. Its gravitational realization is formulated
by means of a tetrad and a connection. This section establishes the first-order
geometry and specifies the type of the map between it and the TPK state space
\@.

\subsection{Discrete Gauss–Stokes identity and incidence transport to the boundary}

Let \(K\) be a cellular complex, let \(W\) be an abelian
group, and let \(\Omega\in C^1(K;W)\) be a cochain. We denote by
\(\delta:C^1(K;W)\to C^2(K;W)\) the cellular coboundary and by \(\langle\,\cdot\,,\,\cdot\,\rangle\) the pairing
between cochains with values in \(W\) and integral cellular chains.

\begin{theorem}[Universal Stokes Cell Identity]
For every \(2\)-cadena finite \(z\in C_2(K;\mathbb Z)\),
\[
 \boxed{\langle\delta\Omega,z\rangle
 =\langle\Omega,\partial z\rangle.}
\]
\end{theorem}
\begin{proof}
This is the identity that defines the coboundary as the adjoint operator of the cellular
boundary. The finiteness of \(z\) ensures that both pairings are finite
sums in \(W\).
\end{proof}

\begin{corollary}[Gauss--Stokes on a cellular pseudomanifold]
Let \(K\) be a finite, compact, regular two-dimensional cellular pseudomanifold,
coherently oriented and possibly with boundary. If \(z_K\) is the sum of
its \(2\)-cells with the coherent orientation, then
\[
 \sum_{f\in K^{(2)}}(\delta\Omega)(f)
 =\sum_e[\partial z_K:e]\,\Omega(e).
\]
In particular, regularity allows the uncancelled incidences of
\(\partial z_K\) to be identified with the edges of the geometric boundary.
\end{corollary}
\begin{proof}
The theorem is applied to \(z_K\). Coherent orientation makes each
interior edge appear twice in \(\partial z_K\) with opposite coefficients;
the uncancelled incidences form the boundary chain.
\end{proof}

The non-Abelian version is not obtained merely by replacing sums with products.
In this work it is used only when transports to a common
base, an ordering of the holonomies, and the corresponding conjugation have been fixed. Outside
that subdomain, no non-Abelian Stokes identity is asserted.

\subsection{Area capacity of the qutrit cell: \(81\log 3\) and collective incidence transport}

\begin{theorem}[Minimum cut of a discrete cube]
In the cubic graph \(N\times N\times N\) with unit capacity, every cut
separating two opposite faces has capacity at least \(N^2\), and one exists
with capacity \(N^2\).
\end{theorem}
\begin{proof}
There are \(N^2\) edge-disjoint paths crossing the cube in the direction
normal to the two faces. Every cut must intersect them, so its capacity is
at least \(N^2\). A transverse layer contains exactly \(N^2\) edges and
attains the bound.
\end{proof}

For \(N=9\), the cut equals 81 and
\[
 9^3=729=27^2.
\]
The last equality permits six trits to be regrouped as a volume \(9^3\) or a
grid of cardinality \(27^2\). It is a bijection of sets; it is not an
isomorphism between
\((\ZZ/9\ZZ)^3\) and \((\ZZ/27\ZZ)^2\), whose exponents are distinct.

\subsection{Tetrad, connection, and \(2\) defects}

Let \(e^I\) be a tetrad and \(\omega^{IJ}=-\omega^{JI}\) a Lorentz
connection. Define
\[
 T_{\rm tor}^I=D_\omega e^I
 =\dd e^I+\omega^I{}_J\wedge e^J,
\]
\[
 F^{IJ}=\dd\omega^{IJ}+\omega^I{}_K\wedge\omega^{KJ},
 \qquad
 g=\eta_{IJ}e^I\otimes e^J.
\]
The decomposition
\[
 \omega^{IJ}=\mathring\omega^{IJ}(e)+K^{IJ}
\]
separates the Levi--Civita connection and the contorsion. Then
\[
 T_{\rm tor}^I=K^I{}_J\wedge e^J.
\]
Curvature and torsion are distinct defects: \(F\) measures closure of
vector transport; \(T_{\rm tor}\), closure of the tetrad. The
\(T\) memory of the Ambivalence Principle is likewise not geometric torsion: only a
realization functor can relate them.

\subsection{1st-order action}

Let \(M\) be a smooth oriented four-dimensional manifold. If
\(\partial M\ne\varnothing\), we assume compactly supported variations in the
interior or, equivalently for the variational problem considered, a
boundary term and boundary conditions that annul the boundary
contribution. When \(\Psi\) contains spinorial fields, we also assume a
spin structure and fields compatible with it. Let
\[
 \Sigma^{IJ}=e^I\wedge e^J,
 \qquad
 (\star X)^{IJ}=\frac12\epsilon^{IJ}{}_{KL}X^{KL},
 \qquad
 \kappa=\frac{8\pi G}{c^4},
 \qquad
 \gamma\in\mathbb R\setminus\{0\}.
\]
The Palatini--Cartan--Holst action is
\[
\boxed{
 S[e,\omega,\Psi]
 =\frac1{2\kappa c}\int_M
 \Sigma_{IJ}\wedge\left(\star F^{IJ}+\frac1\gamma F^{IJ}\right)
 -\frac\Lambda{\kappa c}\int_M\operatorname{vol}_e
 +S_{\rm m}[e,\omega,\Psi].}
\]
The variation of this action determines the torsional corrections
\cite{Sciama1964,Kibble1961,Hehl1976,Holst1996}.

We define the spin current by
\[
 \delta_\omega S_{\rm m}
 =-\frac12\int_M\delta\omega^{IJ}\wedge\sigma_{IJ}.
\]

\begin{theorem}[Cartan--Holst Equation]
Under the preceding geometric and boundary hypotheses, variation with respect
to \(\omega\) gives
\[
 D_\omega\!\left[(\star+\gamma^{-1}I)\Sigma^{IJ}\right]
 =\kappa c\,\sigma^{IJ},
\]
with the global normalization fixed by the preceding definition of
\(\sigma^{IJ}\).
\end{theorem}
\begin{proof}
One uses \(\delta F^{IJ}=D_\omega\delta\omega^{IJ}\), integrates by parts, and
sets the coefficient of the arbitrary variation
\(\delta\omega^{IJ}\) equal to zero. With the sign convention fixed in the definition of
\(\sigma^{IJ}\), the matter term moves to the right-hand side with a
positive sign. The factor \(1/(2\kappa c)\) produces the coefficient
\(\kappa c\). The cosmological term does not depend on \(\omega\), and the
boundary conditions eliminate the term generated by integration by
parts.
\end{proof}

In signature lorentziana \(\star^2=-I\) on bivectores and, for
\(\gamma\in\RR\setminus\{0\}\),
\[
 (\star+\gamma^{-1}I)^{-1}
 =\frac{\gamma^2}{1+\gamma^2}(-\star+\gamma^{-1}I).
\]

\begin{corollary}[Torsional decoupling]
For a co-tetrada no degenerada and
\(\gamma\in\mathbb R\setminus\{0\}\),
\[
 \sigma^{IJ}=0\Longrightarrow T_{\rm tor}^I=0
 \Longrightarrow\omega=\mathring\omega(e).
\]
\end{corollary}

To make the normalization explicit, we define the normalized sources by
the action through
\[
 \delta_g S_{\rm m}
 =-\frac12\int_M\operatorname{vol}_e\,
 \mathcal T^{S}_{\mu\nu}\,\delta g^{\mu\nu}.
\]
When the tetrad coordinates have the dimension of length, the conventional
physical energy--momentum tensor is
\(T^{\rm phys}_{\mu\nu}:=c\,\mathcal T^{S}_{\mu\nu}\). We likewise denote by
\(\mathcal U^{S,\rm spin}\) the contribution normalized by the
action and by \(U^{\rm phys,\rm spin}:=c\,\mathcal U^{S,\rm spin}\) its physical
convention. Upon algebraically eliminating the connection, the metric
equation then assumes the two equivalent forms
\[
 G_{\mu\nu}+\Lambda g_{\mu\nu}
 =\kappa c\,
 \bigl(\mathcal T_{\mu\nu}^{S,\rm LC}
       +\mathcal U_{\mu\nu}^{S,\rm spin}\bigr)
 =\kappa\,
 \bigl(T_{\mu\nu}^{\rm phys,LC}
       +U_{\mu\nu}^{\rm phys,spin}\bigr),
\]
where the spin contribution is quadratic in the corresponding current
and depends on the matter coupling. The coefficient of that contribution is not
fixed without specifying the fermionic action and the duality
conventions.

\subsection{Bianchi Identity with Torsion and Flow Balance in the Cartan Connection}

For every connection,
\[
 D_\omega T_{\rm tor}^I=F^I{}_J\wedge e^J,
 \qquad
 D_\omega F^{IJ}=0.
\]
After removing torsion,
\[
\nabla^\mu(T_{\mu\nu}^{\rm phys,LC}
          +U_{\mu\nu}^{\rm phys,spin})=0.
\]
These identities are mandatory consistency checks: every HMT
gravitational term must arise from a compatible action or satisfy the corresponding
Noether balance.

\subsection{Map of realization and specialization of the Holst parameter}

The insertion defined in \cref{mdg:chap:barbero} takes
\[
 \gamma=\Gamma_{6\to8}.
\]
Once this substitution has been made, all the preceding variational consequences
are theorems of the specialized action. The physical selection of
this action by HMT transport is expressed through a realization map
\[
 \mathfrak R_{\rm grav}:\mathcal X_{\rm TPK}
 \dashrightarrow\{(e,\omega,\Psi)\}
\]
endowed with three properties: intertwining TPK transport with that of
\(\omega\), identifying the holonomy classes, and furnishing the scales of
\(G\) and \(\Lambda\). These properties fix the precise domain of a
gravitational realization of the formalism.

For a nondegenerate co-tetrad, nonzero real \(\gamma\), and matter without spin current, the torsional sector reduces to the first-order formulation of general relativity with the same cosmological constant, subject to the already declared boundary conditions. In a spin fluid, the scaling \(s^2\propto a^{-6}\) can produce an Einstein--Cartan bounce; obtaining its density and initial conditions from the TPK state is a question distinct from the variational consistency of the action.
